\section{Application: a Bitcoin data-availability layer for an optimistic rollup}
\label{sec:rollup-da}

Section~\ref{sec:bitcoin-annex} treated the annex as a place to commit an
arbitrary structured file. This section sketches a specific consumer of
that mechanism: an optimistic rollup that uses a Bitcoin transaction annex
as its data-availability (DA) layer, with FlatMerkle giving watchtowers a
compact way to point at exactly the step they dispute. This is an
architectural proposal, not a protocol this paper implements or
benchmarks; the composition below is intended to be checked, not taken on
faith, and several of its pieces are explicitly flagged as open design
choices rather than settled specifications.

\subsection{Setup: the state-transition commitment}
\label{sec:rollup-setup}

A rollup sequencer periodically batches a sequence of $m$ transactions
$\mathit{tx}_1, \dots, \mathit{tx}_m$ applied to its ledger, starting from a
known initial state $\mathit{state}_0$ (the previous batch's final,
already-committed state), and producing $\mathit{state}_1, \dots,
\mathit{state}_m$, one per transaction. The transactions themselves must be
published \emph{in plain} as part of the committed data, not reduced to a
hash: a data-availability layer that only commits to $H(\mathit{tx}_i)$
without publishing $\mathit{tx}_i$ itself gives nobody the bytes they would
need to recompute $\mathit{state}_i$ and dispute it, defeating the point of
using an annex for DA at all. And since each step's starting state is just
the previous step's ending state, $\mathit{state}_{i-1}$ needs no separate
storage of its own: the sequencer publishes exactly $\mathit{state}_0$
followed by the pairs $(\mathit{tx}_1, \mathit{state}_1), \dots,
(\mathit{tx}_m, \mathit{state}_m)$, one initial value plus one
transaction-and-resulting-state per step, as a single FlatMerkle structured
file (Section~\ref{sec:construction}) inside the annex of a Bitcoin
transaction it signs, using the block-alignment convention of
Section~\ref{sec:annex-alignment}. Anyone who downloads the full annex can
replay $\mathit{state}_0 \to \mathit{state}_1 \to \cdots \to \mathit{state}_m$
directly and confirm the sequencer's claimed final state follows from its
claimed initial state -- an ordinary $O(m)$ check that needs nothing from
this paper. FlatMerkle's contribution is what it adds on top: an
$O(\log N)$ way to point at, and prove the exact bytes of, one disputed
step without redistributing the whole batch.

Packing this sequence into FlatMerkle's fixed-width elements $D[a]$ needs
an explicit framing rule, because each $\mathit{tx}_i$ has variable length
and occupies as many consecutive elements as its encoding requires, while
each $\mathit{state}_i$, a fixed-width digest, occupies exactly one. A
dispute proof must show not only the bytes of some elements but that those
elements \emph{are} $\mathit{tx}_i$ and its neighbouring states. We
therefore let the file begin with a \emph{step table}: element $0$ holds
$\mathit{state}_0$, and the next $\lceil m \cdot e / b \rceil$ elements
hold, for each step $i$, the index of the first element of
$\mathit{tx}_i$ and its length in elements, encoded in a fixed $e$-bit
entry; $\mathit{state}_i$ is the element immediately after
$\mathit{tx}_i$. A table whose entries overlap, leave gaps, run past $N$,
or are not in increasing order makes the batch invalid in the same way an
incorrect transition does, and is disputed the same way. A watchtower
disputing step $i \ge 2$ then proves, in one batched multi-element proof
(Section~\ref{sec:multiproof}), the table element(s) holding entries
$i-1$ and $i$, the elements of $\mathit{tx}_i$, and the elements holding
$\mathit{state}_{i-1}$ and $\mathit{state}_i$; for $i = 1$, the entry for
step $1$, $\mathit{tx}_1$, $\mathit{state}_1$, and element $0$. The proof
has size $O(\ell_i + \log N)$, where $\ell_i$ is the number of elements
$\mathit{tx}_i$ occupies.

Embedding this same commitment in an \texttt{OP\_RETURN} output instead of
the annex raises a separate question -- how a structured file of this kind
maps onto \texttt{OP\_RETURN}'s own size and structuring conventions --
that this paper does not address; everything here is specific to the
annex.

\subsection{Two independent watchtower checks}
\label{sec:rollup-checks}

A set of watchtowers, each holding the full published annex, run two
independent checks against every batch.

\begin{enumerate}
\item \textbf{Canonicity of the layout itself.} A watchtower runs
  $\textsc{ScanForViolations}$ (Algorithm~\ref{alg:scan}, $O(N)$ where $N$
  is the number of elements in the batch's structured file) against the
  published annex. If the file is not canonically built, binding
  (Theorem~\ref{thm:binding}) still holds -- each index has at most one
  accepted value -- but that value need not be the bytes physically
  present at the corresponding position (Remark~\ref{rem:layout-forgery}).
  A light client following the batch through inclusion proofs and a full
  node replaying the raw annex bytes would then see two different batches.
  A batch that fails this check has published something that cannot be
  relied on as a
  faithful, unambiguous state-transition record at all, independent of
  whether any individual transition it contains happens to be correct,
  which this paper treats as sufficient grounds on its own for the
  surrounding rollup protocol to slash the sequencer, hand off sequencing
  to a new party, or force an exit to Layer 1 -- the specific choice among
  these is a protocol-design question this paper does not resolve.
\item \textbf{Correctness of an individual state transition.} For a
  canonical batch, a watchtower checks each step $\mathit{state}_{i-1}
  \to \mathit{state}_i$ against its own independent execution of the
  rollup's deterministic state-transition function. If it finds one step
  $i$ where the published $\mathit{state}_i$ does not match, it can isolate
  exactly that step with a short FlatMerkle inclusion proof
  (Section~\ref{sec:algorithm}, batched per Section~\ref{sec:multiproof} as
  described above) against the batch's committed digest, without needing
  any other watchtower to re-derive or re-publish the rest of the batch.
\end{enumerate}

\subsection{Enforcing the checks on Bitcoin}
\label{sec:rollup-enforce}

Both checks above are things any off-chain party can already do by
downloading the annex; the open problem is making a \emph{failed} check
economically consequential -- getting Bitcoin's own consensus, which
cannot evaluate SHA-256 compression-function chaining or an arbitrary
state-transition function directly in Script, to adjudicate the dispute
and pay out a slashing bond. Two current, concrete mechanisms exist for
this, and either fits here depending on how the rollup structures its own
state-transition check: \textbf{BitVMX}, a virtual CPU that runs an
arbitrary program (RISC-V or MIPS) and resolves a dispute by an $n$-ary
search over a hash chain of the execution trace down to one instruction,
adjudicated by a small, fixed piece of Script \citep{lerner2024bitvmx};
and \textbf{BitVM3}, a newer, garbled-circuit-based protocol purpose-built
to verify a single SNARK proof, where a challenger evaluates a garbled
verifier circuit off-chain and, on a false claim, posts a compact,
constant-size fraud proof on-chain \citep{linuswoll2026bitvm3}. We use
BitVMX below since it fits an arbitrary re-execution check directly; a
rollup that instead produces one succinct SNARK proof per batch could use
BitVM3 for the same role at lower on-chain dispute cost.

Concretely, a watchtower that found a violation of either kind opens a
BitVMX dispute whose program takes as input the FlatMerkle inclusion proof
for the disputed index (or, for a canonicity violation, the fraud proof
produced by $\textsc{BuildFraudProof}$, Algorithm~\ref{alg:build-fraud}), the
revealed step data, $\mathrm{sha\_annex}$, and enough of the Bitcoin
transaction and its Taproot signature to reconstruct $\mathit{SigMsg}$
(Section~\ref{sec:annex-authentication}). The disputed program itself
checks, in order: that the inclusion (or fraud) proof folds correctly to
the claimed $\mathrm{sha\_annex}$ under $\mathrm{finalize\_annex}$
(Section~\ref{sec:annex-alignment}); that the supplied signature validates
over $\mathit{SigMsg}$ for the claimed $\mathrm{sha\_annex}$; and, for a
transition dispute, that re-applying the state-transition function to the
revealed $\mathit{state}_{i-1}$ and $\mathit{tx}_i$ does \emph{not} yield
the revealed $\mathit{state}_i$ (for a canonicity dispute, that
$\textsc{VerifyFraudProof}$ accepts the supplied fraud proof). If BitVMX's
bisection game resolves in the watchtower's favor, the pre-arranged
slashing path pays out. Authenticating a signature against $\mathit{SigMsg}$
this way establishes only the first of the two sub-properties identified in
Section~\ref{sec:annex-authentication} -- that a validly signed witness
variant with this $\mathrm{sha\_annex}$ exists, not that it is the specific
variant Bitcoin confirmed. Chaining successive batch-commitment
transactions (each spending an output of the previous one) does not close
that gap: a spend references the previous transaction's \texttt{txid},
which excludes witness data, so it does not fix which annex variant was
mined. For the dispute, however, sub-property~1 is enough, provided the
slashing condition covers \emph{every} witness variant the sequencer
signs, not only the confirmed one: an invalid (non-canonical, or
transition-violating) annex that carries a valid sequencer signature is
slashable whether or not it was mined, so the dispute program never needs
to learn which variant Bitcoin confirmed, and the watchtowers, who run full
nodes, dispute the mined one. A light client that wants to learn the
rollup's \emph{confirmed} state from a FlatMerkle proof still faces the
trilemma of Section~\ref{sec:annex-authentication}.

The linkage between an externally-signed Bitcoin transaction and a value
consumed inside a BitVMX program is itself the subject of ESSPI
\citep{lerner2025esspi}, which replaces a direct in-Script Schnorr
verification -- reported there as producing a roughly 1.1~GB Script when
attempted directly -- with a Winternitz-committed hash of the signed data,
checked by a secondary BitVMX instance that the commitment, the signature,
and the claimed data are mutually consistent. ESSPI's own construction is
built around a transaction signed by a counterparty \emph{as part of} the
linking protocol, and its authors report no formal security analysis of
the composed protocol and flag exposure to Bitcoin's transaction-replacement
behavior; whether it applies directly to authenticating an
already-existing, independently signed batch-commitment transaction of the
kind described here, or needs adaptation, is not something either ESSPI or
this paper settles.

\subsection{Scope}
\label{sec:rollup-scope}

This section is a design sketch composing three separately published,
separately analyzed mechanisms -- FlatMerkle's inclusion and fraud proofs,
BitVMX's dispute game, and ESSPI's signature-linkage technique -- and none
of the composition itself has been implemented, benchmarked, or given a
security proof here. Left unaddressed: the concrete slashing, sequencer-
succession, and exit-mode game the surrounding rollup protocol would need;
collateral sizing and the round/cost trade-offs BitVMX itself exposes as
configurable \citep{lerner2024bitvmx}; the encoding of the step table's
entries; and, as already noted, embedding a structured file of this kind in
\texttt{OP\_RETURN} rather than the annex.
